\section{Solutions for Chapter~\ref{ch:unitarity}}

\paragraph{Exercise~\ref{ex:unitarity-damping}.}
\phantomsection
\label{sol:ex:unitarity-damping}
Set $a=-2\log r>0$. The continuous function
$f(x)=(x+1)e^{-ax}$ has derivative
$e^{-ax}[1-a(x+1)]$. Thus
\[
 \sup_{n\ge0}(n+1)r^{2n}\le
 \begin{cases}
 1,&a\ge1,\\
 e^{a-1}/a,&0<a<1.
 \end{cases}
\]
Choosing an integer nearest $1/a-1$ also shows that the
discrete supremum diverges like $1/(ea)$ as $a\downarrow0$.
The estimate is finite for every proper annulus and loses
uniformity as the damping annulus degenerates.



\paragraph{Exercise~\ref{ex:unitarity-projector}.}
\phantomsection
\label{sol:ex:unitarity-projector}
Because $P$ preserves the common core, $1-P$ preserves it too,
and $T_0(1-P)=(1-P)T_0$. Apply the same graph-limit
argument to both projections. Their ranges then give a
direct-sum decomposition of the \emph{closed operator and its
domain}. The adjoint of this direct sum is the direct sum of
the adjoints, proving reduction of $T^*$. Only proving
invariance of one subspace would not by itself give this
direct-sum statement.



\paragraph{Exercise~\ref{ex:unitarity-counting}.}
\phantomsection
\label{sol:ex:unitarity-counting}
Full faithfulness identifies each image endomorphism algebra
with $\C$, so the image of a simple is simple in the target
semisimple category. It also preserves the vanishing of
Hom spaces between distinct simples, making the twelve
images pairwise nonisomorphic. They exhaust the twelve target
simples, and additivity exhausts finite direct sums.
A tensorator would be a coherent natural isomorphism
$F(X)\otimes F(Y)\to F(X\otimes Y)$.
Counting objects and morphism dimensions neither supplies
these maps nor verifies their pentagon and braiding compatibilities.



\paragraph{Exercise~\ref{ex:unitarity-deck}.}
\phantomsection
\label{sol:ex:unitarity-deck}
A full rotation on the $h=1/3$ module is
$e^{2\pi i/3}$, while on the integer-graded vacuum VOA it is
one. If $v\in U$ is homogeneous, every mode changes energy
by the integer $\wt(v)-n-1$. Hence every mode commutes
with $Z=e^{2\pi iL_0}$, and
$BZ^mY(v,z)Z^{-m}A=BY(v,z)A$ on the test domain.
Bounded extension preserves this equality. The paired
cancellation does not make the individual charged-sector
phase equal to one.



\paragraph{Exercise~\ref{ex:unitarity-end}.}
\phantomsection
\label{sol:ex:unitarity-end}
For four mutually inequivalent summands, Schur's lemma gives
$\C^4$, of dimension four. If $Y\cong Y^*$, the $Y$
multiplicity space is two-dimensional, giving
$\C\oplus\C\oplus M_2(\C)$, of dimension six and
noncommutative. The canonical extension argument translates
these dimensions into the rank-four and rank-six alternatives.
It is the repeated simple, rather than four displayed summands,
that produces a matrix block.



